\begin{frame}{Integration of Short Term Learning Methods for Image Retrieval by Reciprocal Rank Fusion}
\scriptsize
\vspace{-0.8em}

\begin{table}[h!]
\centering
\renewcommand{\arraystretch}{1.5} % increases row height slightly more
\setlength{\tabcolsep}{2.5pt}     % reduces horizontal padding a bit more

\begin{tabular}{|p{2.6cm}|p{4.3cm}|p{4.3cm}|}
\hline
\textbf{Citation} & \textbf{Summary} & \textbf{Advantages \& Disadvantages} \\
\hline
\tiny\textit{[8] B. Bagheri, M. Pourmahyabadi, and H. Nezamabadi-pour, 2013 8th Iranian Conf. Mach. Vis. Image Process. (MVIP), pp. 366–369, 2013.  \newline
DOI: 10.1109/IranianMVIP.2013.6780012.}
&
Proposes a \textbf{CBIR system} combining two \textbf{STL} methods - \textbf{J\_R-N} and \textbf{SVM Active Learning} - using \textbf{RRF}. Extracts color, texture, and shape features, applies both \textbf{STL} methods independently, and fuses ranked lists using \textbf{RRF}. Tested on the Corel dataset with 10,000 images across 82 semantic groups, achieving 73\% precision after 4 feedback iterations.
&
\textbf{Advantages:}
\begin{itemize}\itemsep0.25em
\item 10\% improvement over J\_R-N and 6\% over \textbf{SVM} individually.
\item Combines complementary strengths of classification and similarity-based approaches.
\item Simple, unsupervised fusion method with straightforward scoring.
\end{itemize}

\textbf{Disadvantages:}
\begin{itemize}\itemsep0.25em
\item Limited to basic feature extraction.
\item Computational cost from multiple method executions and \textbf{RRF}.
\item Tested only on Corel dataset; generalizability uncertain.
\end{itemize} \\
\hline
\end{tabular}
\end{table}
\end{frame}
